% TEMPLATE-MILC-v3.tex - plantilla maestra para todas las guias MILC v3
% Ver scripts/generadores/build-guia-milc.py para la lista completa de
% claves esperadas. El script valida que no queden placeholders sin
% reemplazar antes de escribir el archivo final.

\documentclass[11pt,letterpaper]{article}
\usepackage[margin=1.75cm]{geometry}
\usepackage[table]{xcolor}
\usepackage{fontspec}
\usepackage[spanish]{babel}
\usepackage{tikz}
\usepackage[most]{tcolorbox}
\usepackage{tabularx}
\usepackage{array}
\usepackage{fancyhdr}
\usepackage{titlesec}
\usepackage{ragged2e}
\usepackage{enumitem}
\usepackage{microtype}

\setmainfont{Helvetica Neue}
\setsansfont{Helvetica Neue}
\setlength{\parindent}{0pt}
\setlength{\parskip}{6pt}
\hyphenpenalty=200
\exhyphenpenalty=200
\pretolerance=400
\tolerance=2000
\setlength{\emergencystretch}{1em}
\renewcommand{\arraystretch}{1.25}
\setlist{leftmargin=5mm,itemsep=1mm,topsep=1mm}
\newcolumntype{Y}{>{\RaggedRight\arraybackslash}X}

\definecolor{milcMagenta}{HTML}{E6007E}
\definecolor{milcVino}{HTML}{5A0038}
\definecolor{milcTurquesa}{HTML}{0093A5}
\definecolor{milcVerde}{HTML}{58A923}
\definecolor{milcMostaza}{HTML}{E5B400}
\definecolor{milcOcre}{HTML}{B86600}
\definecolor{milcNegro}{HTML}{241816}
\definecolor{milcGris}{HTML}{F7F7F4}
\definecolor{milcLinea}{HTML}{E8E2DD}
\definecolor{milcGradePrimary}{HTML}{7C3AED}
\definecolor{milcGradeDark}{HTML}{4C1D95}
\definecolor{milcGradeSoft}{HTML}{EDE9FE}
\definecolor{milcGradeAccent}{HTML}{CFE7FF}
\definecolor{milcGradeText}{HTML}{FFFFFF}
\color{milcNegro}

\pagestyle{fancy}
\fancyhf{}
\lhead{\small Tecnología e Informática · Grado 9}
\rhead{\small Guía 11 · MILC v3}
\cfoot{\small\thepage}
\renewcommand{\headrulewidth}{0pt}

\newtcolorbox{titlebox}[1]{
  enhanced,colback=#1,colframe=#1,arc=7mm,boxrule=0pt,
  left=7mm,right=7mm,top=3mm,bottom=3mm,
  before skip=8pt,after skip=8pt,
  fontupper=\sffamily\bfseries\Large\color{white}
}

\newtcolorbox{softbox}[3]{
  enhanced,colback=#2,colframe=#1,borderline west={4pt}{0pt}{#1},
  arc=2mm,boxrule=.4pt,left=4mm,right=4mm,top=3mm,bottom=3mm,
  before skip=6pt,after skip=8pt,title={#3},coltitle=white,
  colbacktitle=#1,fonttitle=\sffamily\bfseries,fontupper=\RaggedRight,
  attach boxed title to top left={xshift=3mm,yshift=-2mm},
  boxed title style={arc=2mm, boxrule=0pt}
}

\newtcolorbox{drawbox}[2]{
  enhanced,colback=white,colframe=#1,arc=4mm,boxrule=1pt,
  left=5mm,right=5mm,top=4mm,bottom=4mm,before skip=8pt,after skip=8pt,
  title={#2},coltitle=white,colbacktitle=#1,fonttitle=\sffamily\bfseries,
  fontupper=\RaggedRight,
  attach boxed title to top left={xshift=4mm,yshift=-2mm},
  boxed title style={arc=3mm, boxrule=0pt}
}

\newtcolorbox{infoband}[1]{
  enhanced,colback=#1!8,colframe=#1!50,arc=2mm,boxrule=.5pt,
  left=4mm,right=4mm,top=2mm,bottom=2mm,before skip=4pt,after skip=4pt,
  fontupper=\small\RaggedRight
}

% Puente narrativo entre secciones (contrato editorial Conéctate).
% Da continuidad: "Acabas de X. Ahora vas a Y porque Z."
% Visualmente: banda izquierda en color del grado, texto en itálica pequeña.
\newtcolorbox{puentebox}{
  enhanced,colback=milcGradeSoft,colframe=milcGradeDark,
  borderline west={3pt}{0pt}{milcGradeDark},
  arc=0mm,boxrule=0pt,left=5mm,right=4mm,top=2.5mm,bottom=2.5mm,
  before skip=4pt,after skip=8pt,
  fontupper=\itshape\small\color{milcNegro}\RaggedRight
}
\newcommand{\puente}[1]{%
  \begin{puentebox}%
    \textcolor{milcGradeDark}{\textbf{\textrightarrow}}\hspace{2mm}#1%
  \end{puentebox}%
}

\newcommand{\linead}[1]{\noindent\color{milcLinea}\rule{#1}{0.5pt}\color{milcNegro}}
\newcommand{\respuesta}[1]{\par\vspace{2mm}\foreach \n in {1,...,#1}{\noindent\color{milcLinea}\rule{\linewidth}{0.5pt}\par\vspace{5mm}}\color{milcNegro}}
\newcommand{\checkbox}{\textcolor{milcGradeDark}{\fbox{\phantom{x}}}\hspace{5pt}}

\newcommand{\phasecard}[4]{
\begin{tcolorbox}[enhanced,colback=#3,colframe=#2,arc=3mm,boxrule=.5pt,height=3.05cm,valign=center,halign=center,left=1.5mm,right=1.5mm,top=2mm,bottom=2mm]
{\sffamily\bfseries\fontsize{22}{24}\selectfont\textcolor{#2}{#1}}\par
{\sffamily\bfseries\small #4}
\end{tcolorbox}
}

\begin{document}
\thispagestyle{empty}
\begin{tikzpicture}[remember picture,overlay]
  \fill[white] (current page.south west) rectangle (current page.north east);
  \path[fill=milcGradePrimary,rounded corners=16mm]
    ([xshift=.55cm,yshift=-.55cm]current page.north west) rectangle
    ([xshift=-.55cm,yshift=3.65cm]current page.south east);

  \node[anchor=north west,text=milcGradeText] at ([xshift=2.0cm,yshift=-1.85cm]current page.north west)
    {\sffamily\bfseries\fontsize{14}{16}\selectfont GRADO};
  \node[anchor=north west,text=milcGradeText,opacity=.82] at ([xshift=2.0cm,yshift=-2.75cm]current page.north west)
    {\sffamily\bfseries\fontsize{9}{11}\selectfont SERIE GUÍAS · TECNOLOGÍA E INFORMÁTICA};

  \begin{scope}[shift={([xshift=-3.05cm,yshift=-2.55cm]current page.north east)}]
    \fill[white,opacity=.80] (-.62,-.52) rectangle (.62,.52);
    \draw[milcGradeText,line width=1pt] (-.62,-.52) rectangle (.62,.52);
    \fill[milcGradeDark] (-.42,-.38) rectangle (-.22,.06);
    \fill[milcGradeAccent] (-.10,-.38) rectangle (.10,.24);
    \fill[milcGradePrimary!60!black] (.22,-.38) rectangle (.42,.38);
  \end{scope}

  \node[anchor=west,text=milcGradeText] at ([xshift=1.9cm,yshift=-12.85cm]current page.north west)
    {\sffamily\bfseries\fontsize{124}{124}\selectfont 9};
  \draw[milcGradeText,line width=1.7pt]
    ([xshift=4.52cm,yshift=-11.68cm]current page.north west) circle (.13cm);

  \node[anchor=west,text=milcGradeText,text width=8.7cm,align=left] at ([xshift=10.35cm,yshift=-11.6cm]current page.north west)
    {
     {\sffamily\bfseries\fontsize{19}{22}\selectfont ¿Qué es diseñar? ---\\orden, jerarquía,\\intención}\\[4mm]
     {\sffamily\bfseries\fontsize{23}{26}\selectfont Guía 11-9-TIC}\\[2mm]
     {\sffamily\fontsize{13}{16}\selectfont Período 2 · Diseño editorial digital}\\[5mm]
     {\sffamily\bfseries\fontsize{11}{13}\selectfont Institución Educativa Sor María Juliana}\\[1mm]
     {\sffamily\fontsize{10}{12}\selectfont Autor: Dr. Álvaro Cárdenas Orozco}
    };

  \path[fill=milcGradeSoft,rounded corners=7mm]
    ([xshift=1.35cm,yshift=.85cm]current page.south west) rectangle
    ([xshift=-1.35cm,yshift=4.25cm]current page.south east);
  \draw[milcGradeDark,line width=.65pt,rounded corners=7mm]
    ([xshift=1.35cm,yshift=.85cm]current page.south west) rectangle
    ([xshift=-1.35cm,yshift=4.25cm]current page.south east);
  \node[anchor=center,text=milcNegro,text width=17.0cm,align=center] at ([yshift=2.55cm]current page.south)
    {\sffamily\fontsize{10}{12}\selectfont Metodología MILC · Apertura ancestral · Pensamiento computacional · Triángulo Dussel-Estoicismo-Floridi\\
    \textcolor{milcGradeDark}{\bfseries Producto:} Análisis comparado de 3 productos editoriales reales (libro, revista, periódico o volante) identificando decisiones de orden, jerarquía e intención del diseñador};
\end{tikzpicture}
\mbox{}
\newpage

\section*{Apertura: ¿Qué es diseñar? --- orden, jerarquía, intención}

\begin{softbox}{milcGradeDark}{milcGradeSoft}{Saber ancestral}
Diseñar no nació con Figma. Los \textbf{pasquines populares} del Valle, los \textbf{almanaques campesinos}, los \textbf{pliegos parroquiales} y las \textbf{cartillas escolares} de los años 50 eran piezas de diseño con propósito: jerarquía clara, tipografía que se podía leer en la plaza, orden de información que respondía a una necesidad concreta de la comunidad. Antes del diseño profesional hubo \emph{oficio editorial} ---  carpintero del texto, sastre de la página, panadero de la información.
\end{softbox}

\begin{softbox}{milcTurquesa}{milcGris}{Saber contemporáneo}
Hoy ``diseñar'' suele asociarse a Figma, Canva o Adobe. Pero el principio es el mismo desde el pasquín: \textbf{decidir qué importa más y comunicarlo con orden}. El diseño editorial digital hereda 500 años de tipografía, jerarquía y composición. Tu próxima publicación en redes sociales también es diseño editorial --- aunque no la pienses así.
\end{softbox}

\begin{softbox}{milcMostaza}{milcGris}{Pregunta-puente}
\textit{¿Qué tienen en común un pasquín del Valle hace 100 años y una historia de Instagram bien hecha hoy? ¿Y en qué se diferencian profundamente?}
\end{softbox}

\begin{softbox}{milcVerde}{milcGris}{Saber y saber hacer}
Al terminar podrás: (1) \textbf{identificar} 3 decisiones de diseño en productos editoriales reales (libro, revista, periódico); (2) \textbf{analizar} cómo \emph{orden, jerarquía e intención} sostienen la legibilidad; (3) \textbf{explicar} la diferencia entre diseño profesional y maquetación apresurada; (4) \textbf{crear} un análisis comparado de 3 productos editoriales con tus propios criterios.
\end{softbox}

\small
\begin{tabularx}{\textwidth}{>{\bfseries}p{3.0cm}Y}
\rowcolor{milcGradePrimary!12}Autor & Dr. Álvaro Cárdenas Orozco\\
DBA & Producción de piezas editoriales digitales con criterio estético y ético (MEN, T\&I)\\
Referentes & Tipografía colombiana · Diseño centrado en accesibilidad · Floridi\\
Producto final & Análisis comparado de 3 productos editoriales reales (libro, revista, periódico o volante) identificando decisiones de orden, jerarquía e intención del diseñador\\
\end{tabularx}
\normalsize

\puente{Antes de empezar, mira el viaje completo del periodo 2. En 10 sesiones vas a construir una revista digital de 8 páginas. Hoy es la apertura: entrenar el ojo a ver decisiones de diseño que otros toman sin que tú lo notes.}

\begin{titlebox}{milcGradeDark}
Ruta MILC: así vamos a aprender
\end{titlebox}
\begin{tabularx}{\textwidth}{>{\centering\arraybackslash}X>{\centering\arraybackslash}X>{\centering\arraybackslash}X>{\centering\arraybackslash}X}
\phasecard{1}{milcVerde}{milcGris}{Escuta\\[-1mm]\small Observo, escucho y pregunto.} &
\phasecard{2}{milcTurquesa}{milcGris}{Sistematización\\[-1mm]\small Pensamiento computacional.} &
\phasecard{3}{milcMagenta}{milcGris}{Praxis\\[-1mm]\small Construyo y aplico.} &
\phasecard{4}{milcVino}{milcGris}{Evaluación\\[-1mm]\small Reflexión liberadora.}
\end{tabularx}


\puente{Antes de teorizar, vamos a leer 3 productos editoriales reales. Detrás de cada uno hay decisiones específicas --- de tamaño, color, orden. Aprender a verlas es el primer paso del oficio.}

\begin{titlebox}{milcVerde}
Fase 1 · Escuta
\end{titlebox}

\begin{softbox}{milcVerde}{milcGris}{Escena para escuchar}
\textbf{Actividad 1 · IDENTIFICA --- 3 productos editoriales reales} (15 min · individual). Sin moverte mucho, busca \textbf{3 productos editoriales} cerca tuyo: (a) un \emph{libro} (texto, ensayo, novela), (b) una \emph{revista} (impresa o digital), (c) un \emph{periódico o volante}. Para cada uno: ¿qué se ve primero al abrir? ¿qué viene después? ¿qué es lo más pequeño?
\end{softbox}

\checkbox Tengo a mano 3 productos editoriales reales (libro, revista, periódico/volante).\\
\checkbox Identifiqué qué se ve primero, segundo, tercero en cada uno.\\
\checkbox Reconocí 1 decisión de diseño en cada producto (tamaño, color, espacio).

\begin{infoband}{milcVerde}
\textbf{Lo que va al cuaderno (Actividad 1 · IDENTIFICA).} Título: «Actividad 1 --- 3 productos editoriales bajo la lupa». \textbf{Formato:} tabla de 4 columnas (Producto~|~Qué se ve primero~|~Qué viene después~|~Decisión de diseño identificada), 3 filas. \textbf{Sabes que terminaste cuando} para cada producto puedes nombrar al menos UNA decisión específica del diseñador.
\end{infoband}

\puente{Ya viste decisiones reales. Ahora vamos a entender los \textbf{3 pilares} de todo buen diseño editorial: orden, jerarquía e intención. Distinguirlos con criterio te permite leer cualquier pieza editorial, vieja o nueva.}

\begin{titlebox}{milcTurquesa}
Fase 2 · Sistematización con pensamiento computacional
\end{titlebox}

\begin{softbox}{milcTurquesa}{milcGris}{Cuatro pilares aplicados al tema}
Todo diseño editorial fuerte se sostiene en 3 pilares. Vamos a descomponerlos con los pilares del pensamiento computacional para que sepas qué buscar en cualquier pieza.
\end{softbox}

\small
\begin{tabularx}{\textwidth}{>{\bfseries\RaggedRight\arraybackslash}p{3.6cm}Y}
\rowcolor{milcTurquesa!90}\color{white}Pilar & \color{white}\bfseries Aplicación al tema\\
1. Descomposición & El diseño editorial se descompone en 3 pilares: (1) \textbf{orden} (qué va antes, qué después --- secuencia narrativa), (2) \textbf{jerarquía} (qué es más importante, qué es secundario --- tamaño, peso, color), (3) \textbf{intención} (para qué sirve la pieza --- informar, persuadir, entretener, recordar).\\
2. Reconocimiento de patrones & Todo diseño fuerte sigue el patrón \textbf{``decisión visible, no accidente''}: en una pieza bien diseñada, cada tamaño, cada color, cada espacio responde a una decisión. En una pieza apresurada, todo se ve igual ---  el ojo no sabe por dónde empezar y la lectura se vuelve trabajo.\\
3. Abstracción & Lo esencial del diseño cabe en una pregunta: \emph{¿qué quiero que el lector lea primero, qué después, y qué no necesita leer?}. Sin esa pregunta, diseñar es decorar.\\
4. Algoritmo conceptual & Pasos para analizar cualquier pieza editorial: (1) ¿cuál es su intención principal?; (2) ¿qué orden propone el diseñador (de arriba abajo, de izquierda a derecha, otra)?; (3) ¿cómo establece jerarquía (tamaño, peso, color, espacio)?; (4) ¿qué decisión te parece más fuerte y cuál más débil?.\\
\end{tabularx}
\normalsize

\begin{drawbox}{milcTurquesa}{Anatomía de una pieza editorial}
\textbf{Intención:} informar, persuadir, entretener, recordar. \\[1mm] \textbf{Orden:} secuencia narrativa que propone (lineal, jerárquica, modular). \\[1mm] \textbf{Jerarquía visual:} tamaños, pesos, colores que distinguen lo importante. \\[1mm] \textbf{Espacio negativo:} aire que respira entre elementos. \\[1mm] \textbf{Tipografía:} familia, peso, alineación. \\[1mm] \textbf{Color:} paleta, contraste, función.
\end{drawbox}

\begin{softbox}{milcMostaza}{milcGris}{Errores comunes y su corrección}
(1) Confundir ``diseñar'' con ``decorar'' --- el diseño tiene función, la decoración es opcional. (2) Pensar que ``si está bonito está bien diseñado'' --- la belleza sin función no convence. (3) Olvidar la intención: pieza con buena tipografía pero sin propósito claro pierde fuerza. (4) Imitar tendencias sin entenderlas --- copiar Figma sin saber por qué se ve así.
\end{softbox}

\begin{infoband}{milcTurquesa}
\textbf{Lo que va al cuaderno (Actividad 2 · EXPLICA).} Título: «Actividad 2 --- 3 pilares del diseño editorial». \textbf{Formato:} 3 fichas, una por pilar (orden, jerarquía, intención), cada una con: definición en tus palabras + 1 ejemplo de uno de los 3 productos analizados en Actividad 1 que lo cumple bien. \textbf{Sabes que terminaste cuando} puedes explicar los 3 pilares sin notas y reconocer cuál es más fuerte en cada producto.
\end{infoband}

\puente{Tienes el método. Toca producir: un análisis comparado escrito de 3 productos editoriales. Sin diseñar todavía --- analizar primero.}

\begin{titlebox}{milcMagenta}
Fase 3 · Praxis con pensamiento computacional
\end{titlebox}

\begin{softbox}{milcMagenta}{milcGris}{Construyendo el producto con los cuatro pilares}
\textbf{Actividad 3 · CREA --- Análisis comparado de 3 productos editoriales} (30 min · individual). Hora de escribir. Vas a hacer un análisis comparado escrito (texto + dibujo/esquema) de 3 productos editoriales reales, evaluando cómo aplican los 3 pilares.
\end{softbox}

\small
\begin{tabularx}{\textwidth}{>{\bfseries\RaggedRight\arraybackslash}p{3.6cm}Y}
\rowcolor{milcMagenta!90}\color{white}Pilar & \color{white}\bfseries Aplicación al producto\\
Descomposición & Tu análisis se descompone en 4 secciones: (1) introducción (qué productos elegiste y por qué), (2) análisis individual (los 3 pilares aplicados a cada uno), (3) comparación cruzada (qué tienen en común, qué los diferencia), (4) conclusión personal (qué decisión de diseño quieres recordar para tu propia revista).\\
Patrones & Sigue el patrón \textbf{``ejemplo específico vence a categoría abstracta''}: no escribas ``el libro tiene buena tipografía''. Escribe ``el libro usa una tipografía con serifa de cuerpo 10pt con interlineado 15pt que facilita la lectura larga''. Lo concreto convence.\\
Abstracción & Cada análisis individual expresa una sola idea principal: cómo logra la pieza su intención. Si tu análisis enumera 7 detalles sin jerarquía propia, perdiste el foco.\\
Algoritmo de armado & Algoritmo: (1) elige 3 productos editoriales distintos; (2) analiza cada uno con los 3 pilares (5 min cada uno); (3) escribe comparación cruzada (5 min); (4) cierra con 1 decisión de diseño que te llevas (5 min); (5) revisa que cada apartado tenga al menos UN ejemplo específico.\\
\end{tabularx}
\normalsize

\begin{drawbox}{milcMagenta}{Checklist antes de cerrar el análisis}
\checkbox 3 productos editoriales distintos analizados (libro + revista + periódico/volante). \\[1mm] \checkbox Cada análisis aplica los 3 pilares (orden, jerarquía, intención). \\[1mm] \checkbox Hay al menos 1 ejemplo específico por análisis (no generalidad). \\[1mm] \checkbox Comparación cruzada identifica al menos 1 común y 1 diferencia. \\[1mm] \checkbox Conclusión personal nombra 1 decisión que vas a recordar para tu propia revista.
\end{drawbox}

\begin{softbox}{milcMagenta}{milcGris}{Plantilla de trabajo}
\textbf{Introducción.} \textit{Qué 3 productos elegí y por qué.} \\[1mm] \textbf{Análisis individual.} \textit{Cada producto con los 3 pilares aplicados.} \\[1mm] \textbf{Comparación cruzada.} \textit{Qué tienen en común, qué los diferencia.} \\[1mm] \textbf{Conclusión personal.} \textit{Una decisión de diseño que voy a recordar para mi revista.}
\end{softbox}

\begin{infoband}{milcMagenta}
\textbf{Lo que va al cuaderno (Actividad 3 · CREA).} Título: «Actividad 3 --- Análisis comparado de 3 productos editoriales». \textbf{Formato:} 4 secciones escritas + dibujo/esquema simple de cada producto. \textbf{Extensión:} 2 páginas de cuaderno. \textbf{Sabes que terminaste cuando} un compañero puede leer tu análisis y entender qué decisión de diseño te llevas.
\end{infoband}

\puente{Lo que sigue es el resumen de qué entregas hoy: análisis comparado escrito. El criterio principal: que identifiques decisiones específicas, no generalidades estéticas.}

\begin{titlebox}{milcMostaza}
Producto de la sesión
\end{titlebox}
\textbf{Análisis comparado de 3 productos editoriales con los 3 pilares aplicados}.
\begin{softbox}{milcMostaza}{milcGris}{Criterios de calidad}
(1) 3 productos editoriales distintos analizados. (2) Los 3 pilares (orden, jerarquía, intención) aplicados a cada uno. (3) Al menos 1 ejemplo específico por producto (no generalidad). (4) Comparación cruzada con común y diferencia. (5) Conclusión personal con decisión accionable para la revista del periodo.
\end{softbox}

\puente{Antes de cerrar, mira el oficio de diseñar desde las cinco dimensiones humanas. Diseñar bien es un acto profundamente político: decides quién puede leer, qué se ve primero, qué se invisibiliza.}

\begin{titlebox}{milcVino}
Fase 4 · Evaluación liberadora — cinco dimensiones
\end{titlebox}

\begin{softbox}{milcVino}{milcGris}{Sentido de la evaluación}
Evaluar no es solo revisar una nota. Es reconocer qué aprendí, qué cambió en mí, qué emociones manejé, cómo me relaciono con la ciudad y la comunidad, y qué le dejaré a quien viene detrás.
\end{softbox}

\textbf{Mi reflexión liberadora en cinco dimensiones.} Completa cada frase iniciada:

\small
\begin{tabularx}{\textwidth}{>{\bfseries\RaggedRight\arraybackslash}p{3.4cm}Y>{\RaggedRight\arraybackslash}p{4.6cm}}
\rowcolor{milcVino}\color{white}Dimensión & \color{white}\bfseries Mi respuesta (frame starter) & \color{white}\bfseries Algo que descubrí\\
1. Personal & Lo que más cambió en mí fue \linead{6.5cm} & \linead{4.2cm}\\
2. Emocional & La emoción más fuerte fue \linead{2.5cm} y la regulé \linead{3cm} & \linead{4.2cm}\\
3. Ciudadana & En democracia esto importa porque \linead{6cm} & \linead{4.2cm}\\
4. Local (Cartago) & En mi barrio sería útil para \linead{6.5cm} & \linead{4.2cm}\\
5. Intergeneracional & A mi abuela/abuelo le diría \linead{6.5cm} & \linead{4.2cm}\\
\end{tabularx}
\normalsize

\begin{infoband}{milcVino}
\textbf{Para la próxima sustentación.} Vuelve a esta tabla en seis meses. Verás que las respuestas cambian — no porque lo aprendido cambie, sino porque tú cambias. La evaluación liberadora es una conversación contigo a través del tiempo.
\end{infoband}


\puente{Y para terminar, tres voces sobre diseñar. Dussel sobre los diseños que excluyen, Marco Aurelio sobre la sencillez bien hecha, Floridi sobre el diseño como organizador de la infoesfera.}

\begin{titlebox}{milcGradeDark}
Triángulo de pensamiento · cierre filosófico
\end{titlebox}

\begin{softbox}{milcVino}{milcGris}{Dussel · lente del nosotros}
\textit{``Diseñar es decidir quién puede leer; toda decisión de diseño excluye a alguien.''}

Una pieza editorial bien diseñada para ojo joven puede ser invisible para ojo anciano. Un sitio web optimizado para celular caro puede excluir a quien usa teléfono básico. Un manual médico en jerga técnica excluye al paciente. Cada decisión de diseño tiene consecuencias políticas: incluye o excluye. La inclusión empieza por reconocer a quién dejaste afuera y diseñar de nuevo.

\textbf{¿A quién excluyen las piezas editoriales que admiro? ¿Y mis propias piezas, a quién dejan afuera sin darme cuenta?}
\end{softbox}
\linead{\linewidth}\\[1mm]
\linead{\linewidth}

\begin{softbox}{milcMostaza}{milcGris}{Estoicismo · Marco Aurelio · cuidado interior}
\textit{``Lo sencillo bien hecho es más difícil que lo complicado mal hecho.''}

La tentación del diseñador novato es ``llenar la página'' con tipografías, colores, decoraciones. La sabiduría estoica del diseño es lo opuesto: \emph{simplicidad rigurosa}. Una página con UNA fuente, dos colores, espacio generoso y jerarquía clara se ve más profesional que una con 5 fuentes y 10 colores. Aprender a quitar es más difícil que aprender a agregar.

\textbf{Cuando diseño algo, ¿agrego o quito? ¿Qué de lo que pongo es necesario y qué es solo ruido visual que me da miedo eliminar?}
\end{softbox}
\linead{\linewidth}\\[1mm]
\linead{\linewidth}

\begin{softbox}{milcTurquesa}{milcGris}{Floridi · lente de la infoesfera}
\textit{``El diseño es el lenguaje silencioso con que la infoesfera se organiza.''}

Cada vez que un diseñador decide tamaño, orden y color, organiza la infoesfera. Las decisiones del que diseñó Twitter (cómo se ven los tweets, qué se destaca) moldean el debate público mundial. El diseño nunca es neutro: \emph{moldea la atención humana a escala masiva}. Aprender a diseñar es aprender a tomar responsabilidad por la atención de otros.

\textbf{¿Qué decisiones de diseño en las plataformas que uso a diario moldean mi atención sin que yo lo note? ¿Cuáles diseñaría diferente?}
\end{softbox}
\linead{\linewidth}\\[1mm]
\linead{\linewidth}

\begin{titlebox}{milcVino}
Compromiso final
\end{titlebox}

\small
\begin{tabularx}{\textwidth}{>{\RaggedRight\arraybackslash}p{3.0cm}YYY}
\rowcolor{milcVino}\color{white}\bfseries Criterio & \color{white}\bfseries Lo logré & \color{white}\bfseries En proceso & \color{white}\bfseries Necesito apoyo\\
Comprensión del tema & Explico ideas centrales con mis palabras. & Reconozco ideas pero falta precisión. & Necesito repasar conceptos base.\\
Pensamiento computacional & Apliqué los 4 pilares al diseñar mi producto. & Apliqué algunos pilares. & Necesito releer la fase 2.\\
Aplicación y producto & Mi producto cumple los criterios y se puede revisar. & Producto funcional con ajustes pendientes. & Aún en construcción.\\
Reflexión 5 dimensiones & Respondí cada dimensión con honestidad. & Respondí algunas con detalle. & Mis respuestas son muy generales.\\
Triángulo de pensamiento & Conecté las 3 voces con mi trabajo real. & Conecté al menos una. & No relaciono las voces con mi proceso.\\
\end{tabularx}
\normalsize

\textbf{Hoy entendí que:}\\[1mm]
\linead{\linewidth}\\[1mm]
\linead{\linewidth}

\textbf{Mi compromiso para la próxima sesión es:}\\[1mm]
\linead{\linewidth}\\[1mm]
\linead{\linewidth}

\begin{softbox}{milcMostaza}{milcGris}{Cita de despedida · Marco Aurelio}
\textit{``El obstáculo en el camino se vuelve el camino. Lo que estorba al camino se vuelve camino.''}\\[1mm]
Cada error de hoy, bien observado, se convierte en pieza del camino que pisarás mañana.
\end{softbox}

\small
\begin{tabularx}{\textwidth}{>{\bfseries}p{3.6cm}Y}
\rowcolor{milcGradeDark!12}Nombre completo: & \linead{10cm}\\
Fecha: & \linead{4cm} \quad Curso/grupo: \linead{4cm}\\
Mi firma: & \linead{9cm}\\
\end{tabularx}
\normalsize

\begin{infoband}{milcGradeDark}
\textbf{Para la próxima semana.} Lleva tu cuaderno durante 3 días y, cada vez que veas una pieza editorial (cartel, volante, libro, post de redes), anota UNA decisión de diseño que te llamó la atención (positiva o negativa). Al cabo de 3 días lee la lista. Vas a notar cuántas decisiones pasan invisibles. El próximo lunes traemos esa bitácora --- es insumo de la sesión 2.
\end{infoband}

\end{document}
